\section{Related Work}

Classical equation discovery follows two main approaches. Sparse regression methods like SINDy \cite{brunton2016sindy,messenger2021weak,fasel2022ensemble} and PDE-FIND \cite{rudy2017pdefind} select active terms from predefined libraries via sparsity-promoting optimization. Genetic programming methods—Eureqa \cite{schmidt2009eureqa}, PySR \cite{cranmer2023pysr}, Operon \cite{burlacu2020operon}—evolve expression trees, with physics-informed variants adding domain constraints \cite{zhang2024pigp}. Neural approaches include AI Feynman \cite{udrescu2020aifeynman,udrescu2020aifeynman2,kamienny2022e2e,mundhenk2021seeding,landajuela2022udsr} exploiting physical symmetries, Deep Symbolic Regression \cite{petersen2021dsr} using risk-seeking policy gradients, and transformer models like NeSymReS \cite{biggio2021nesymres} and SymFormer \cite{vastl2022symformer}. KBASS \cite{long2023kbass} combines Bayesian spike-and-slab priors with kernel methods. SRBench \cite{lacava2021srbench} and Feynman equations \cite{udrescu2020aifeynman} provide standard benchmarks. Recent LLM-based methods leverage pre-trained scientific knowledge: LLM-SR \cite{shojaee2024llmsr,boiko2023coscientist,ma2024sga,zhang2024chemllm} uses evolutionary refinement on Python-represented equations, LLM4ED \cite{du2024llm4ed} alternates self-improvement and evolutionary phases, FunSearch \cite{romera2024funsearch} achieved novel mathematical discoveries, LaSR \cite{grayeli2024lasr} builds reusable concept libraries, and ICSR \cite{merler2024icsr} applies in-context learning. However, these methods use static prompts, unable to adapt instructions based on search progress. Prompt optimization methods include APE \cite{zhou2023ape} using bandit selection, OPRO \cite{yang2024opro} with LLMs as meta-optimizers, and evolutionary approaches like EvoPrompt \cite{guo2024evoprompt} and PromptBreeder \cite{fernando2024promptbreeder}. Bayesian prompt optimization has also been explored \cite{sabbatella2024promptbo}. Yet existing methods seek a single optimal prompt, whereas we argue \textit{different prompts are optimal at different stages}. Our work introduces closed-loop instruction optimization via Bayesian Optimization, dynamically selecting strategies based on numerical feedback.

